\section{Related Work}

\paragraph{Bidirectional Video Generation Models.}
Video generation has advanced rapidly in recent years, with modern approaches mostly adopting the paradigms of denoising diffusion. Video diffusion has been explored in both pixel space~\citep{ho2022imagen, singer2022make} and latent space~\citep{blattmann2023align, blattmann2023stable}, with architectures evolving from early Space–Time U-Nets~\citep{blattmann2023stable, hong2022cogvideo} to more recent DiT-based designs~\citep{peebles2023scalable, gupta2024photorealistic}. Significant industrial investment has driven the development of large video diffusion models, leading to several multi-billion parameter models, including open-source models such as Wan~\citep{wan2025wan} and Hunyuan~\citep{kong2024hunyuanvideo}, and closed-source models such as Sora~\citep{openai_sora}, Movie-Gen~\citep{polyak2024movie}, and Seaweed~\citep{seawead2025seaweed}. Notably, these models operate as bidirectional video diffusion models, as they have access to both past and future frames during denoising. While this bidirectional context enables high-quality synthesis for offline generation, it is incompatible with the causality that is necessitated in real-time streaming video generation.


\paragraph{Autoregressive Video Generation Models.}
To enable long video generation, several studies have extended the generation paradigm from bidirectional to autoregressive, which naturally supports gradual rollout over extended time horizons. Autoregressive models are typically trained with next-token prediction objectives and generate spatiotemporal tokens sequentially at inference time~\citep{bruce2024genie, kondratyuk2023videopoet, wang2024loong, weissenborn2019scaling, yan2021videogpt}. More recently, a separate line of research combines autoregressive modeling with denoising diffusion~\citep{chen2024diffusion, gu2025long, guo2025long, jin2024pyramidal, li2024arlon, liu2024mardini, weng2024art, yin2025slow, zhang2025test, zhang2025packing, huang2025self, henschel2025streamingt2v}, where frames are generated one-by-one in an outer loop and each frame is gradually denoised in an inner loop. Within this family, Rolling Diffusion~\citep{ruhe2024rolling} and its variants~\citep{kim2024fifo, teng2025magi, sun2025ar, xie2025progressive, chen2025skyreels, teng2025magi} merge the outer and inner loops: the diffusion model jointly denoises multiple frames at progressively increasing noise levels. However, these methods mostly suffer from exposure bias and error accumulation when generating long videos. Another line of research addresses error accumulation with planning generation~\citep{long2024videostudio, zhao2024moviedreamer, hu2024storyagent, xie2024dreamfactory, zhang2025packing, bansal2024talc, yang2024synchronized, xiang2025macro}, which predicts distant future frames first and then interpolates the intermediate frames. While effective for reducing drift, it breaks the strict sequential order required for real-time streaming. In contrast, our work enables much longer streaming video generation with minimal error accumulation while addressing exposure bias.



\paragraph{Concurrent and Closed-Source Work.} 
Two concurrent works are also devoted to the streaming generation of long videos.
% 
Specifically, StreamDiT~\citep{kodaira2025streamdit} adopts the FIFO-style denoising~\citep{kim2024fifo} for streaming video generation. It modifies the pretrained model architecture by introducing micro-steps and window attention, necessitating extensive additional pretraining with large-scale data and computation. In contrast, our method keeps the pretrained model architecture unchanged, can be trained efficiently in only 3,000 steps, and does not require any video data.
APT2~\citep{lin2025autoregressive} instead explores adversarial distillation~\citep{lin2025diffusion} for streaming video generation. It denoises videos block-by-block and involves multiple costly post-training stages, including diffusion adaptation, consistency distillation, adversarial training, and long-video training. APT2 is trained on one-minute-long videos, whereas ours is trained only on 5-second clips, yet can extend to multi-minute sequences during inference.
Note that both StreamDiT and APT2 are closed-source and trained based on internal video diffusion models (Movie-Gen~\citep{polyak2024movie} and Seaweed~\citep{seawead2025seaweed}), 
while our model is trained on public datasets and relies on an open-source model (\ie, Wan2.1~\citep{wan2025wan}) as its foundation.
% They also have far larger model sizes (8B and 4B parameters, respectively) compared to ours (1.3B). 
